\Article{39}
The value of $v$, which in general is a function of $r$, $u$ and $\lambda$, becomes a function of $u$ and $\lambda$ alone on the Earth's surface, and
\[
-\frac{\mathrm dv}{R\mathrm du},\qquad -\frac{\mathrm dv}{R\sin u\,\mathrm d\lambda}
\]
are the horizontal components of the magnetic force produced there by the galvanic currents, directed northward and westward respectively. It is therefore evident that the remarkable propositions cited above in Articles 15 and 16 apply here as well. For the third component, the vertical magnetic force, the situation differs somewhat depending on whether the agents are located above the Earth or in its interior. To determine the vertical force arising from the former, one must first take $v$ as a function of all three variables $r$, $u$ and $\lambda$, differentiate it with respect to $r$ and then substitute $r=R$. For the inner region $S'$, to which the Earth's surface belongs, however, $v$ can only be expanded in a series of ascending powers of $r$. If we put
\[
\frac vR=p^0+\frac rR\cdot p'+\frac{rr}{RR}\cdot p''+\frac{r^3}{R^3}\cdot p'''+\text{etc.}
\]
then $p^0$ is a constant, namely the value of $\frac vR$ at the centre of the
\SourceLeaf{176}
Earth. By contrast, $p'$, $p''$, $p'''$, etc. are functions of $u$ and $\lambda$ that satisfy the same partial differential equations as $P'$, $P''$, $P'''$, etc. above. It follows, in much the same way as in Article 20 above, that knowing the value of $v$ at every point on the Earth's surface is sufficient to derive the general expression valid throughout $S'$. Knowing just the horizontal forces on the Earth's surface is sufficient to obtain that value apart from a constant part, or equivalently, to determine the coefficients $p'$, $p''$, $p'''$, etc. On the same surface, however, the vertical force does not have the value
\[
=2p'+3p''+4p'''+\text{etc.}
\]
(which it would have if the forces were produced from within the Earth). Its value is
\[
=-p'-2p''-3p'''-\text{etc.}
\]

Now our numerical elements (Article 26), determined on the assumption of the first formula, already give a very satisfactory representation of the phenomena as a whole, whereas those phenomena would be entirely incompatible with the second formula. The untenability of the hypothesis that places the causes of terrestrial magnetism in the space outside the Earth must therefore be regarded as established.

\Article{40}
Even so, this cannot yet be taken to disprove conclusively the possibility that a \textit{part} of the Earth's magnetic force, even if only a comparatively very small part, is produced from above. A much fuller and much more precise knowledge of the phenomena will shed light on this important point of the theory in the future. On the assumption of mixed causes, let the symbols $V$, $P^0$, $P'$, $P''$, etc., $v$, $p^0$, $p'$, $p''$ have the same meanings as above, with the former referring to the causes acting from within the Earth and the latter to the causes acting from the space outside it. If we further set
\[
V+v=W,\quad P^0+p^0=\Pi^0,\quad P'+p'=\Pi',\quad P''+p''=\Pi''\quad\text{etc.}
\]
\SourceLeaf{177}
then on the Earth's surface we have
\[
\frac WR=\Pi^0+\Pi'+\Pi''\quad\text{etc.}
\]
where $\Pi^{\mathrm n}$ satisfies the same partial differential equation as $P^{\mathrm n}$ (Article 18), and the two components of the horizontal magnetic force acting there are expressed by
\[
-\frac{\mathrm dW}{R\mathrm du},\qquad-\frac{\mathrm dW}{R\sin u\,\mathrm d\lambda}
\]
The propositions cited in Articles 15 and 16 therefore still hold here. From the horizontal forces alone one can determine the quantities $\Pi'$, $\Pi''$, $\Pi'''$, etc., but they alone tell us nothing about whether mixed causes are present. If, however, the vertical force is considered separately and put in the form
\[
Q^0+Q'+Q''+Q'''+\text{etc.}
\]
so that $Q^{\mathrm n}$ satisfies the partial differential equation mentioned above. then we have
\[
\begin{gathered}
Q^0=P^0,\quad Q'=2P'-p',\quad Q''=3P''-2p'',\quad Q'''=4P'''-3p'''
\end{gathered}
\]
etc., and consequently
\[
\begin{aligned}
3P'&=\Pi'+Q',&3p'&=2\Pi'-Q'\\
5P''&=2\Pi''+Q'',&5p''&=3\Pi''-Q''\\
7P'''&=3\Pi'''+Q''',&7p'''&=4\Pi'''-Q'''\quad\text{etc.}
\end{aligned}
\]

Combining the horizontal forces with the vertical force thus gives us the means to separate $W$ into its parts $V$ and $v$ and to find out whether a perceptible value must be assigned to the latter. The observations leave only the constant part of $v$, namely $p^0$, wholly undetermined. The reason is evident from Article 38.

From this interesting point of view too, it therefore appears important to consider the horizontal magnetic force on its own, and we see in this a further reason to heed the recommendations made above in Article 21.

\SourceLeaf{178}
\Article{41}
There will probably be a shortage of adequate data for the investigation outlined in the preceding article for a long time to come. It is worth noting, however, that the variations in the magnetic force. as they occur simultaneously at different points on the Earth's surface, lend themselves to a very similar treatment, for which it may be possible to collect the minimum data needed much sooner. This applies both to the regular changes that vary with the time of day and of year and to the irregular changes. There may be room here for a few general remarks about these future investigations.

Once the simultaneous changes observed at each place have been put in the form of changes in the components of the magnetic force, $\Delta X$, $\Delta Y$, $\Delta Z$, the first task is to examine whether the changes in the two horizontal components behave in accordance with our theory, according to which $-\Delta X$ and $-\sin u.\Delta Y$ must be the values of the partial derivatives of a function of $u$ and $\lambda$ with respect to these variables. If they do, we conclude that the causes are either actual galvanic currents or at least act in the same way as these currents or as separated magnetic fluids. In the opposite case it would be established that the causes cannot be galvanic currents. We see that knowledge of such changes in the horizontal force alone can provide highly important insights, given sufficient precision, number and geographical spread. If the simultaneous changes in the vertical force are also known, then, \textit{assuming the former case}, the method of the preceding article will show whether the causes are located above or below the Earth's surface. Indeed, since they are probably situated in a layer rather thin compared with the whole body of the Earth, it will then be possible to determine at least approximately how they are distributed.

As for the second case mentioned above as possible, I believe I may regard it as unlikely for the regular changes in the Earth's magnetic force that depend on the time of day and of year. For the irregular changes that vary over short intervals, however, I would hardly venture to make a conjecture
\SourceLeaf{179}
in this regard at present. If they originate in large movements of electricity above the atmosphere, these would hardly belong in the category of galvanic currents. For although everything points to a galvanic current being electricity in motion, such motion constitutes a galvanic current only when it forms a circuit that closes on itself. Since only under this condition is the repeatedly mentioned substitution of separated magnetic fluids for the galvanic current permissible, our relations between the components would cease to hold under that hypothesis. That is, the second case would actually occur. Establishing this important circumstance with certainty would itself be of great interest. Moreover, even then, with sufficiently widespread and reliable observations, it would be within our reach to trace the locations and behaviour of such movements.
